% Phase 1 Report -- Dataset Collection & Tokenizer Construction
% Model H: Hindi. Model L: Bhojpuri.
% Compile with: xelatex phase1_report.tex  (run twice for the ToC/refs)
\documentclass[11pt]{article}

\usepackage[margin=1in]{geometry}
\usepackage{amsmath}
\usepackage{fontspec}
\usepackage{graphicx}
\usepackage{booktabs}
\usepackage{longtable}
\usepackage{array}
\usepackage{xcolor}
\usepackage{caption}
\usepackage{amssymb}
\PassOptionsToPackage{hyphens}{url}
\usepackage[colorlinks=true,linkcolor=blue!50!black,urlcolor=blue!50!black,citecolor=blue!50!black]{hyperref}
\emergencystretch=3em
\sloppy

% Devanagari-capable font for sample tokenizations (Windows ships Nirmala UI)
\newfontfamily\devanagarifont{Nirmala UI}

\newcommand{\dn}[1]{{\devanagarifont #1}}
\newcommand{\pass}{\checkmark}
\newcommand{\fig}[3]{%
  \begin{figure}[htbp]
    \centering
    \includegraphics[width=0.88\textwidth]{figures/#1}
    \caption{#2}
    \label{fig:#3}
  \end{figure}
}

\title{\textbf{Phase 1 Report: Dataset Collection and Tokenizer Construction} \\
\large Model H: Hindi \quad\textbar\quad Model L: Bhojpuri}
\author{Swaraj Kumar \\ \small 2025202016}
\date{\today}

\begin{document}
\maketitle

\begin{abstract}
\noindent
This report covers Phase 1 of the Language Models and Agents individual
project: dataset collection, cleaning, splitting, and tokenizer
construction for two independent, from-scratch decoder-only Transformer
language models. Model H (higher-resource) is \textbf{Hindi}; Model L
(lower-resource, from the assignment's allowed list) is \textbf{Bhojpuri}.
Every number in this report is read directly from a file produced by the
project's own pipeline code -- \texttt{clean\_stats.json},
\texttt{hindi\_pipeline\_manifest.json}, \texttt{validate\_final.log}, or
the trained tokenizer model files re-measured directly -- never estimated.
Both languages clear the assignment's hard 20\% manual-collection floor;
Hindi exceeds the 500M-token planning target (101.2\%) while Bhojpuri's
documented shortfall (16.8\%) is explicitly permitted by the assignment
for the low-resource model. A parallel, exploratory attempt to build
Model L in \textbf{Nepali} instead is also documented in
Section~\ref{sec:nepali}, along with why Bhojpuri was kept as the final
choice.
\end{abstract}

\tableofcontents
\newpage

% =====================================================================
\section{Overview}
% =====================================================================

The assignment requires two independent, from-scratch decoder-only
Transformer language models in two different Indian languages: Model H
from a higher-resource language, Model L from a fixed list of
lower-resource languages (Assamese, Bhojpuri, Bodo, Dogri, Konkani,
Maithili, Manipuri, Mizo, Nepali, Sindhi). Both must be built with fully
separate data, tokenizer, vocabulary, and (from Phase 2) weights -- no
sharing, no multilingual training.

\textbf{Model H is Hindi.} Hindi is one of the highest-resource Indian
languages for NLP: Wikipedia, curated corpora like AI4Bharat's Sangraha,
and web-crawled corpora like Fineweb-2 all have large Hindi partitions.
This makes it realistic to reach the assignment's $\sim$500M-token
planning target almost entirely from public sources, while still
requiring genuine manual collection to clear the 20\% floor.

\textbf{Model L is Bhojpuri}, chosen from the assignment's allowed
low-resource list. Bhojpuri is spoken by tens of millions of people but
has very little curated NLP data: most available corpora are small,
scattered, or web-crawled with heavy noise. This makes Bhojpuri a genuine
test of the assignment's manual-collection requirement -- public sources
alone do not come close to 500M tokens, so the manual share matters more
here by necessity, not just by rule.

Both languages are Devanagari-script, which lets the cleaning and
tokenization code share logic (\texttt{scripts/clean\_lib.py}) while
keeping every byte of data, every trained model file, and every decision
fully separate per language.

% =====================================================================
\section{Data Collection}
% =====================================================================

\subsection{What counts as ``manual'' under the course ruling}

A source counts as \textbf{manual} only if it required work beyond
pointing a downloader at an existing public dataset: a self-run scrape, a
personally extracted Wikipedia XML dump, OCR from a physical or scanned
source, or similarly hands-on collection. A dataset that is just
\emph{labeled} as Wikipedia-derived on Hugging Face does not count --
even though the underlying text originates from Wikipedia, using the
HF-hosted mirror is a download, not manual collection.

Concretely: \texttt{hi\_wikipedia\_dump} is the \textbf{self-extracted}
Hindi Wikipedia XML dump, run through WikiExtractor by this project -- it
counts as manual. \texttt{wikimedia/wikipedia:20231101.hi}, the
HF-mirrored parquet of the same underlying Wikipedia content, does
\textbf{not} count as manual, and using both would double-count Wikipedia
in two different lanes with two different \texttt{source\_type} labels
for the same underlying text. It was moved to
\texttt{hindi/data/raw/\_quarantined\_hf\_hiwiki\_mirror/} and excluded
from cleaning entirely.

\subsection{Hindi (Model H) raw sources}

\begin{table}[htbp]
\centering
\begin{tabular}{lrr}
\toprule
 & Docs & Characters \\
\midrule
Manual & 229{,}011 & 540{,}367{,}035 \\
Downloaded & 544{,}595 & 1{,}299{,}990{,}112 \\
\midrule
\textbf{Total raw} & \textbf{773{,}606} & \textbf{1{,}840{,}357{,}147} \\
\bottomrule
\end{tabular}
\caption{Hindi raw corpus, by lane.}
\end{table}

\begin{table}[htbp]
\centering
\small
\begin{tabular}{lrrp{5cm}}
\toprule
Source & Docs & Characters & Collection method \\
\midrule
\texttt{hi\_wikipedia\_dump} & 160{,}591 & 234{,}096{,}553 & self-extracted Wikipedia XML dump (WikiExtractor) \\
\texttt{gadya\_kosh} & 21{,}481 & 197{,}490{,}954 & scrape + clean (prose wiki) \\
\texttt{patrika\_hindi} & 16{,}251 & 73{,}759{,}826 & scrape + clean (newspaper archive) \\
\texttt{kavita\_kosh} & 30{,}621 & 23{,}610{,}585 & scrape + clean (poetry wiki) \\
archive.org (misc, 67 items) & 67 & 11{,}409{,}117 & OCR/text extraction \\
\bottomrule
\end{tabular}
\caption{Hindi manual sources.}
\end{table}

\begin{table}[htbp]
\centering
\begin{tabular}{p{6cm}rr}
\toprule
Source & Docs & Characters \\
\midrule
HuggingFaceFW/fineweb-2:hin\_Deva & 217{,}460 & 500{,}000{,}289 \\
ai4bharat/sangraha:unverified/hin & 191{,}987 & 500{,}000{,}700 \\
ai4bharat/sangraha:verified/hin:web & 118{,}785 & 230{,}941{,}855 \\
ai4bharat/sangraha:verified/hin:pdf & 16{,}120 & 67{,}642{,}616 \\
ai4bharat/sangraha:verified/hin:speech & 243 & 1{,}404{,}652 \\
\bottomrule
\end{tabular}
\caption{Hindi downloaded sources.}
\end{table}

\textbf{Excluded:} \texttt{wikimedia/wikipedia:20231101.hi} (163{,}091
docs / 253.3M chars) -- the HF-hosted Hindi Wikipedia mirror, quarantined
to avoid double-counting against the self-extracted dump above, and
because it would not count as manual regardless.

\emph{Note on the raw total:} \texttt{clean\_stats.json}'s own
\texttt{chars\_in} (what \texttt{clean\_hindi.py} measured as its input)
reads 1{,}840{,}356{,}619 -- 528 characters (0.00003\%) less than the sum
above. Both figures are real, independently-collected measurements (one at
ingestion/tagging time, one at cleaning time); the gap is negligible and
not investigated further given its scale.

\subsection{Bhojpuri (Model L) raw sources}

\begin{table}[htbp]
\centering
\begin{tabular}{lrr}
\toprule
 & Docs & Characters \\
\midrule
Manual & 47{,}958 & 90{,}781{,}624 \\
Downloaded & 560{,}305 & 227{,}825{,}102 \\
\midrule
\textbf{Total raw} & \textbf{608{,}263} & \textbf{318{,}606{,}726} \\
\bottomrule
\end{tabular}
\caption{Bhojpuri raw corpus, by lane.}
\end{table}

\begin{table}[htbp]
\centering
\small
\begin{tabular}{lrrp{5cm}}
\toprule
Source & Docs & Characters & Collection method \\
\midrule
\texttt{khabar\_bhojpuri} & 22{,}175 & 39{,}169{,}105 & scrape + clean \\
\texttt{anjoria} & 13{,}556 & 21{,}750{,}949 & scrape + clean \\
\texttt{bho\_wikipedia\_dump} & 8{,}539 & 6{,}101{,}826 & self-extracted Wikipedia XML dump \\
\texttt{jogira} & 3{,}085 & 6{,}595{,}231 & scrape + clean \\
archive.org (misc, 88 items) & 88 & 16{,}210{,}212 & OCR/text extraction \\
\texttt{harsh\_scraped\_manual} & 290 & 724{,}804 & prior personal scrape \\
\texttt{sirijan} & 225 & 229{,}497 & scrape + clean \\
\bottomrule
\end{tabular}
\caption{Bhojpuri manual sources. All four dedicated crawl sources
(\texttt{khabar\_bhojpuri}, \texttt{anjoria}, \texttt{jogira},
\texttt{sirijan}) were recrawled uncapped and confirmed exhausted
(0--7 new docs on the final pass).}
\end{table}

\begin{table}[htbp]
\centering
\begin{tabular}{p{6cm}rr}
\toprule
Source (largest four of seventeen) & Docs & Characters \\
\midrule
Satyam810/BhojpuriCorpus:fish\_food & 215{,}813 & 22{,}728{,}391 \\
harsh\_indiccorp (personal prior collection) & 69{,}433 & 92{,}803{,}579 \\
Satyam810/BhojpuriCorpus:jnu\_bhltr & 79{,}884 & 7{,}099{,}650 \\
Satyam810/BhojpuriCorpus:hplt & 17{,}778 & 45{,}689{,}298 \\
\bottomrule
\end{tabular}
\caption{Bhojpuri downloaded sources (largest four; full 17-source table in
\texttt{dataset\_statistics.md}).}
\end{table}

\begin{table}[htbp]
\centering
\begin{tabular}{lp{11cm}}
\toprule
Source & Reason excluded \\
\midrule
bhojpuriplanet.com & domain expired \\
bhojpuri.blog & DNS-unregistered \\
CulturaX, sil-ai/bloom-lm & gated behind an authenticated HF token \\
AI4Bharat Sangraha, Glot500 & zero Bhojpuri rows in the current release \\
MADLAD-400 & no \texttt{bho}/\texttt{bho\_Deva} config \\
CC-100 & legacy \texttt{datasets}-library script format, no longer supported \\
\bottomrule
\end{tabular}
\caption{Bhojpuri sources evaluated and excluded, with reasons.}
\end{table}

\textbf{Open item, documented rather than hidden:} the BHLTR corpus
(\texttt{bhltr\_monolingual}/\texttt{bhltr\_monolingual\_v02}, git-cloned
from \url{https://github.com/shashwatup9k/BHLTR}, CC-BY-NC-SA-4.0) likely
overlaps with \texttt{Satyam810/BhojpuriCorpus:jnu\_bhltr}, which
redistributes the same underlying corpus under a different label. The two
are not deduplicated by source name specifically, but the pipeline's
content-hash dedup (\texttt{dedup\_key}, Section~\ref{sec:cleaning}) still
catches any exact-text overlap between them regardless of label.

\emph{Note on the raw total:} the same kind of small gap exists here --
\texttt{clean\_stats.json}'s \texttt{chars\_in} reads 318{,}528{,}901,
77{,}825 characters (0.024\%) below the 318{,}606{,}726 source-table sum,
for the same reason as Hindi's gap above.

\textbf{A real, measured cross-contamination risk worth naming
explicitly:} \texttt{kavita\_kosh} (a Hindi manual source) is a large
poetry wiki that also hosts smaller Bhojpuri, Awadhi, and Maithili
sections mixed into the same MediaWiki namespace -- the crawl itself is
not language-pure at the point of collection. This is exactly the kind of
risk the language-ID filter (Section~\ref{sec:cleaning}) exists to catch,
and it is why that filter's coarse-marker limitation (documented honestly
below) matters in practice, not just in theory.

% =====================================================================
\section{Cleaning Pipeline}
\label{sec:cleaning}
% =====================================================================

Both languages share one cleaning implementation,
\texttt{scripts/clean\_lib.py}, run separately per language
(\texttt{clean\_hindi.py} / \texttt{clean\_bhojpuri.py}) against that
language's own \texttt{raw/} tree only. The measured stage order,
confirmed below from \texttt{clean\_stats.json}'s own checkpoints (not
assumed from documentation), is:
\begin{quote}\centering
raw $\rightarrow$ malformed-Unicode removal $\rightarrow$
too-short removal $\rightarrow$ language-ID filter $\rightarrow$
exact/near-duplicate removal $\rightarrow$ cleaned
\end{quote}

\textbf{Why this order is verified, not assumed:} \texttt{clean\_stats.json}
records a mid-pipeline checkpoint, \texttt{docs\_interim}, taken after
malformed+short-doc removal but before the language filter. For Hindi:
$773{,}606 - 17 - 63{,}132 = 710{,}457 = \texttt{docs\_interim}$ exactly.
Then $710{,}457 - 3{,}371 - 2{,}983 = 704{,}103 = \texttt{docs\_cleaned}$
(pre-cross-language-correction figure -- see below). The same arithmetic
holds exactly for Bhojpuri.

\textbf{Normalization: NFC, not NFKC.} Unicode NFC (canonical composition)
is used, not NFKC (compatibility composition). NFKC would strip or alter
U+200C (ZWNJ) and U+200D (ZWJ) -- zero-width joiner/non-joiner characters
that look like formatting noise but are \textbf{semantically meaningful}
in Devanagari: they control whether a consonant cluster forms a conjunct
ligature or stays visually separate. Stripping them would silently
corrupt real orthographic distinctions in both Hindi and Bhojpuri text.
NFC preserves them; only genuine zero-width/control characters unrelated
to Devanagari orthography are stripped separately.

\textbf{Latin-token removal, digits kept.} Hindi: 8{,}016{,}102 Latin
tokens removed; Bhojpuri: 1{,}576{,}446. Digits are kept, since numerals
are not script-specific noise the way stray Latin words are.

\textbf{Language filter.} \texttt{langdetect} has \textbf{no \texttt{bho}
label} -- it cannot itself identify Bhojpuri text, and Hindi/Bhojpuri
share vocabulary and script closely enough that this is a real problem.
Quoted directly from each language's \texttt{clean\_stats.json}:

\begin{quote}\small
Hindi (\texttt{bhojpuri\_reject\_caveat}): \textit{``langdetect has no
`bho' label, so it cannot itself catch Bhojpuri text misclassified as
Hindi (they share vocabulary and script).
\texttt{dropped\_language.likely\_bhojpuri} is an active reject filter,
not just a count: any document where Bhojpuri lexical markers
(\dn{बा}/\dn{बाटे}/\dn{बाड़ें}/\dn{एगो}/\dn{गइल}/\dn{भइल}/\dn{कइल}/etc.,
$\geq$2 hits) outnumber Hindi markers
(\dn{है}/\dn{हैं}/\dn{एक}/\dn{गया}/\dn{हुआ}/\dn{किया}/etc.) is dropped
before reaching \texttt{data/cleaned/}. This is a coarse lexical
heuristic, not a language model, so it will not catch every case --
documented as a known limitation, not a solved problem.''}
\end{quote}

\begin{quote}\small
Bhojpuri (\texttt{language\_id\_caveat}): \textit{``langdetect has no
`bho' label; Devanagari-script-ratio was used as a practical proxy and
cannot distinguish Bhojpuri from closely-related Hindi/Nepali text
mislabeled at the source. \texttt{likely\_hindi\_contamination} is a
coarse lexical-marker count among docs that otherwise passed cleaning --
it is \textbf{NOT} a filter, documents counted here are still included in
\texttt{data/cleaned/}.''}
\end{quote}

On the Hindi side, \texttt{likely\_bhojpuri} is an \textbf{active reject
filter} -- 1{,}365 documents dropped. On the Bhojpuri side,
\texttt{likely\_hindi\_contamination} (\textbf{11{,}164} documents) is a
\textbf{measurement only} -- flagged but not removed.

\textbf{This limitation is not hypothetical.}
\texttt{scripts/validate\_phase1.py}'s cross-language check independently
caught 2 documents whose content hash was shared between the Hindi and
Bhojpuri cleaned corpora -- genuinely Bhojpuri news text sourced as
\texttt{hin} by \texttt{ai4bharat/sangraha} while the same text was
correctly labeled \texttt{bho\_Deva} by \texttt{HuggingFaceFW/fineweb-2}.
Both were removed from Hindi's cleaned corpus. This is why Hindi's final
cleaned count is \textbf{704{,}101}, not \texttt{clean\_stats.json}'s own
recorded 704{,}103.

\textbf{Deduplication.} Exact-hash matching on normalized text
(\texttt{dedup\_key}); on a tie between a manual-lane and a
downloaded-lane document, the \textbf{manual copy wins}. Near-duplicate
detection beyond exact-hash matching (e.g. MinHash/LSH) was \textbf{not}
run on the Hindi/Bhojpuri pipelines -- see Section~\ref{sec:limitations}.

\begin{table}[htbp]
\centering
\begin{tabular}{lrr}
\toprule
Reason & Hindi & Bhojpuri \\
\midrule
Malformed Unicode & 17 & 23 \\
Too short ($<$200 chars) & 63{,}132 & 4{,}665 \\
Low Devanagari ratio & 228 & 1{,}501 \\
Likely Bhojpuri (active reject, Hindi only) & 1{,}365 & --- \\
langdetect mismatch & 1{,}778 & --- \\
Exact/near-duplicate & 2{,}983 & 94{,}080 \\
Cross-language duplicate (targeted correction) & 2 & 0 \\
\midrule
\textbf{Total dropped} & \textbf{69{,}505} & \textbf{100{,}269} \\
\textbf{Docs in $\rightarrow$ cleaned} & 773{,}606 $\rightarrow$ 704{,}101 & 608{,}263 $\rightarrow$ 507{,}994 \\
\bottomrule
\end{tabular}
\caption{Measured drop counts, both languages.}
\end{table}

\fig{03_cleaning_funnel.png}{Cleaning funnel: documents remaining after each filter stage, in the pipeline's real, verified order.}{funnel}
\fig{04_drop_reasons.png}{Drop reasons breakdown, ranked, per language.}{dropreasons}

Too-short documents are by far the largest single loss for Hindi
(90.8\% of all drops); for Bhojpuri, exact/near-duplicate dominates
(93.8\%) -- several overlapping \texttt{Satyam810/BhojpuriCorpus}
partitions turned out to share substantial exact-duplicate text once
normalized.

% =====================================================================
\section{Splits}
% =====================================================================

Both languages split at the \textbf{document level}, seed 42, using
deterministic hashing of \texttt{document\_id} so that the same document
can never land in two different splits and no split boundary depends on
processing order.

\begin{table}[htbp]
\centering
\begin{tabular}{lrr}
\toprule
 & Hindi (train/val/test) & Bhojpuri (train/val/test) \\
\midrule
Documents & 675{,}937 / 14{,}082 / 14{,}082 & 487{,}674 / 10{,}160 / 10{,}160 \\
Ratio & $\sim$95.9\% / 2.0\% / 2.0\% & $\sim$96.0\% / 2.0\% / 2.0\% \\
\bottomrule
\end{tabular}
\caption{Train/validation/test split sizes.}
\end{table}

\texttt{scripts/validate\_phase1.py}'s leakage check -- a
\texttt{document\_id} set intersection across all three splits per
language -- reports \textbf{no overlap} for both languages
(\texttt{validate\_final.log}, lines 10--11 and 20--21).

\fig{05_split_sizes.png}{Train/validation/test split sizes, both languages, log scale.}{splits}

% =====================================================================
\section{Tokenizers}
% =====================================================================

\subsection{Why SentencePiece Unigram, not BPE}

Both tokenizers are SentencePiece \textbf{Unigram}, trained separately per
language on that language's own \texttt{data/splits/train/} only.

BPE builds a vocabulary by greedily merging the most frequent adjacent
pair of symbols, one merge at a time, with no mechanism to reconsider an
earlier merge. Unigram instead treats the \emph{segmentation itself} as a
latent variable: it starts from a large candidate set of subword pieces,
then uses EM to iteratively re-estimate piece probabilities and prune the
candidate set down to the target vocabulary size.

This matters for Hindi and Bhojpuri specifically because both are
morphologically rich, agglutinative-suffixing languages layered on top of
Devanagari's conjunct-cluster orthography (a single visual glyph can
correspond to a multi-codepoint sequence forming a consonant conjunct).
BPE's greedy, order-dependent merging has no way to revisit a bad early
merge inside a conjunct cluster or a common suffix chain; Unigram's
probabilistic re-estimation can. This is a structural argument, not a
measured one for this project -- no side-by-side BPE run was performed
here to confirm it empirically (Section~\ref{sec:limitations}).

\textbf{The cost:} Unigram's suffix-array construction stage is
memory-heavy. On this project's 16GB training machine, Hindi's
suffix-array stage was the binding memory constraint, peaking near 5.9GB.

\subsection{Full sweep table}

Vocabulary range: \textbf{5{,}000--10{,}000}, per a course-TA ruling
capping vocab size at this range for a $\sim$25M-parameter model (see
embedding budget below). Every number below was \textbf{re-measured
fresh} for this report by \texttt{report/phase1/make\_plots.py}, loading
each trained \texttt{.model} file directly and encoding the same held-out
sample every existing measurement used: the first 2{,}000 non-empty
documents from \texttt{data/splits/validation/text/}, sorted-filename
order. The fresh numbers match \texttt{hindi\_pipeline\_manifest.json}
and \texttt{bhojpuri\_tokenizer\_sweep.log} exactly to 4+ decimal places.

\begin{table}[htbp]
\centering
\begin{tabular}{rrrrrrr}
\toprule
Vocab & chars/tok & fertility & UNK & byte-fb & coverage & unused ids \\
\midrule
5{,}000 & 3.3230 & 1.4961 & 0.000000 & 0.018301 & 0.9654 & 173 \\
8{,}000 & 3.5622 & 1.3956 & 0.000000 & 0.019613 & 0.9762 & 190 \\
10{,}000 & 3.6601 & 1.3583 & 0.000000 & 0.020213 & 0.9767 & 233 \\
\bottomrule
\end{tabular}
\caption{Hindi tokenizer sweep.}
\end{table}

\begin{table}[htbp]
\centering
\begin{tabular}{rrrrrrr}
\toprule
Vocab & chars/tok & fertility & UNK & byte-fb & coverage & unused ids \\
\midrule
5{,}000 & 3.1314 & 1.6037 & 0.000000 & 0.001649 & 0.9230 & 385 \\
6{,}000 & 3.2394 & 1.5502 & 0.000000 & 0.001706 & 0.9122 & 527 \\
8{,}000 & 3.3900 & 1.4814 & 0.000000 & 0.001785 & 0.8776 & 979 \\
10{,}000 & 3.4935 & 1.4375 & 0.000000 & 0.001840 & 0.8321 & 1{,}679 \\
\bottomrule
\end{tabular}
\caption{Bhojpuri tokenizer sweep.}
\end{table}

\fig{06_vocab_chars_per_token.png}{Vocab sweep: characters per token, both languages.}{cpt}
\fig{07_vocab_fertility.png}{Vocab sweep: fertility (tokens per word), both languages.}{fert}
\fig{08_vocab_byte_fallback.png}{Vocab sweep: byte-fallback rate, both languages.}{bf}

\subsection{Why UNK rate is meaningless here}

UNK rate reads \textbf{exactly 0.0} at every vocab size for both
languages. This is not a sign the vocabulary covers everything -- it is
structural. \texttt{byte\_fallback=True} means any character sequence the
model cannot represent as a normal subword piece falls back to raw-byte
pieces (\texttt{<0xNN>} tokens) instead of the \texttt{<unk>} token. So
UNK never fires; the metric that actually tracks coverage is the
\textbf{byte-fallback rate}. Hindi's byte-fallback rate (1.83\%--2.02\%)
is roughly 11x Bhojpuri's (0.16\%--0.18\%) -- consistent with Hindi's more
heterogeneous source mix.

\subsection{Vocab-size choice: 8,000 for both languages}

At a $\sim$25M-parameter model with tied input/output embeddings, the
embedding table costs $\text{vocab\_size} \times \text{model\_dim}$
parameters. The course-TA ruling ($\sim$49\% of the model at 24k vocab)
corresponds almost exactly to $\text{model\_dim}=512$:
$24{,}000 \times 512 = 12{,}288{,}000 = 49.15\%$ of 25{,}000{,}000.

\begin{table}[htbp]
\centering
\begin{tabular}{rrr}
\toprule
Vocab & Embed params ($d{=}512$) & \% of 25M \\
\midrule
8{,}000 & 4{,}096{,}000 & \textbf{16.38\%} \\
10{,}000 & 5{,}120{,}000 & 20.48\% \\
24{,}000 (ruled out) & 12{,}288{,}000 & 49.15\% \\
\bottomrule
\end{tabular}
\caption{Embedding parameter budget at $d=512$.}
\end{table}

8{,}000 keeps the embedding table a clear minority of the 25M budget
while capturing most of the achievable fertility/utilization gain over
5{,}000. Measured directly: \textbf{8{,}000 produces more total tokens
than 10{,}000 for Hindi} (505{,}852{,}028 vs. 491{,}883{,}387), pushing
Hindi to 101.2\% of the 500M target versus 10{,}000's 98.4\%.

\subsection{Sample tokenizations (real script output, vocab\_size=8,000)}

SentencePiece's word-start marker (the raw output piece \texttt{U+2581},
shown as ``\_'' below purely for font-compatibility in this PDF) precedes
every piece that begins a new whitespace-delimited word; pieces without it
attach to the previous piece with no space.

\textbf{Hindi} (\texttt{hindi\_analyze\_8000\_final.log}):

\begin{quote}\small
\dn{text: अतुल माहेश्वरी के निधन के कई महीनों बाद अमर उजाला में...} \\
pieces: [`\dn{\_अ}', `\dn{तुल}', `\dn{\_माह}', `\dn{ेश्वर}', `\dn{ी}', `\dn{\_के}', `\dn{\_निधन}', `\dn{\_के}', \ldots] \\
ids: [357, 4081, 1847, 1993, 279, 259, 3143, 259, \ldots]

\medskip
\dn{text: नई दिल्ली} \\
pieces: [`\dn{\_नई}', `\dn{\_दिल्ली}', `\texttt{<0x0A>}', `\dn{नर}', `\dn{ेंद्र}', `\dn{\_मोदी}', \ldots] \\
ids: [648, 513, 13, 1796, 2061, 713, \ldots]

\medskip
\dn{text: दिल्ली: अदालत ने सैयद अहमद बुखारी के तर्क को बताया...} \\
pieces: [`\dn{\_दिल्ली}', `:', `\dn{\_अदालत}', `\dn{\_ने}', `\dn{\_सैयद}', `\dn{\_अहमद}', \ldots] \\
ids: [513, 388, 1802, 273, 7550, 3344, \ldots]
\end{quote}

\textbf{Bhojpuri} (\texttt{bhojpuri\_tokenizer\_sweep.log}, vocab\_size=8000):

\begin{quote}\small
\dn{text: कांग्रेस अब आरोप थोपो पार्टी बन कर रह गई है।} \\
pieces: [`\dn{\_कांग्रेस}', `\dn{\_अब}', `\dn{\_आरोप}', `\dn{\_}', `\dn{थो}', `\dn{पो}', `\dn{\_पार्टी}', \ldots] \\
ids: [680, 346, 1011, 263, 4466, 1263, 707, \ldots]

\medskip
\dn{text: तूने जो कौड़ियां भेंट दी हैं वे किसकी दी हुई हैं।} \\
pieces: [`\dn{\_तू}', `\dn{ने}', `\dn{\_जो}', `\dn{\_कौ}', `\dn{ड़िया}', `\dn{ं}', `\dn{\_भेंट}', \ldots] \\
ids: [661, 403, 706, 2714, 3049, 308, 2037, \ldots]

\medskip
\dn{text: कांग्रेसियन के आवे वाला पीढ़ी ले निश्चिंत बा कि ओकर...} \\
pieces: [`\dn{\_कांग्रेस}', `\dn{ियन}', `\dn{\_के}', `\dn{\_आवे}', `\dn{\_वाला}', `\dn{\_पीढ़ी}', \ldots] \\
ids: [680, 494, 259, 544, 320, 3583, \ldots]
\end{quote}

\fig{09_token_frequency_loglog.png}{Token frequency distribution (log-log), vocab\_size=8000, both languages. The Zipf-shaped curve shows a small number of high-frequency pieces (postpositions, punctuation) covering most tokens, with a long tail.}{loglog}

\subsection{Training sample size, stated plainly}

SentencePiece does not train on the entire train split. It uses
\texttt{input\_sentence\_size=2{,}000{,}000} with
\texttt{shuffle\_input\_sentence=True}, reservoir-sampling up to
2{,}000{,}000 sentences from the full train split.

\begin{table}[htbp]
\centering
\begin{tabular}{lrrr}
\toprule
Language & Chars used & Train split total chars & Fraction sampled \\
\midrule
Hindi & $\sim$372.7M & 1{,}692{,}881{,}268 & 22.0\% \\
Bhojpuri & 217{,}511{,}027 & 271{,}000{,}251 & 80.3\% \\
\bottomrule
\end{tabular}
\caption{Fraction of the train split SentencePiece actually saw during training.}
\end{table}

This is SentencePiece's standard, documented mechanism for bounding
training cost -- its own trainer log warns ``Too many sentences are
loaded'' and recommends exactly this option. The held-out evaluation
above is run separately, on the validation split, so it measures the
trained tokenizer's real generalization regardless of how much of train
it saw while fitting. Full-corpus token counts (Section~\ref{sec:results})
are exact and unsampled.

% =====================================================================
\section{Results}
\label{sec:results}
% =====================================================================

\textbf{Hindi:} 505{,}852{,}028 total tokens (train 485{,}280{,}741 +
validation 10{,}484{,}640 + test 10{,}086{,}647) -- \textbf{101.2\% of
the $\sim$500M target}. Manual share: \textbf{28.65\% char-weighted}
(505{,}509{,}566 / 1{,}764{,}596{,}782), \textbf{30.66\% token-weighted}
(155{,}104{,}651 / 505{,}852{,}028).

\textbf{Bhojpuri:} 84{,}217{,}032 total tokens (train 80{,}610{,}590 +
validation 2{,}085{,}234 + test 1{,}521{,}208) -- \textbf{16.8\% of the
$\sim$500M target}. Manual share: \textbf{28.59\% char-weighted}
(80{,}843{,}912 / 282{,}815{,}177), \textbf{29.89\% token-weighted}
(25{,}169{,}871 / 84{,}217{,}032).

\fig{10_token_count_vs_target.png}{Total token count vs. the 500M planning target, both languages.}{target}

\textbf{Bhojpuri's shortfall is expected and spec-permitted, not a gap in
effort.} The assignment explicitly allows the lower-resource model to
fall short of the 500M-token target if public sources cannot reach it,
provided the shortfall is honestly documented (no padding or duplication
was done -- Section~\ref{sec:cleaning}). All four dedicated manual crawl
sources were recrawled uncapped and confirmed exhausted; every
Bhojpuri-labeled HF dataset was evaluated. Published prior Bhojpuri
monolingual resources are, at most, on the order of a 45{,}000-sentence
corpus (JNU BHLTR) and a $\sim$169{,}000-word web corpus -- this project's
282.8M-character cleaned corpus is orders of magnitude larger than
anything previously publicly available for this language.

% =====================================================================
\section{Reproducibility}
% =====================================================================

\begin{table}[htbp]
\centering
\begin{tabular}{ll}
\toprule
Item & Value \\
\midrule
Split seed & 42 (both languages, document-level) \\
Tokenizer determinism & SentencePiece \texttt{shuffle\_input\_sentence=True} \\
Git commit (Hindi pipeline) & \texttt{0b0012ecb3deb08b61a81cc6e989c9bc60289a31} \\
Python & 3.14.0 \\
sentencepiece & 0.2.2 \\
Platform & Windows-11-10.0.26200-SP0 \\
\bottomrule
\end{tabular}
\caption{Reproducibility metadata.}
\end{table}

\textbf{To rerun a stage} (per language, \texttt{<lang>} = \texttt{hindi}
or \texttt{bhojpuri}):

\begin{quote}\small\ttfamily
python <lang>/data/pipeline/download\_<lang>.py \\
python <lang>/data/pipeline/ingest\_manual.py \\
python <lang>/data/pipeline/ingest\_wikipedia\_dump.py \\
python <lang>/data/pipeline/clean\_<lang>.py \\
python <lang>/data/pipeline/split\_<lang>.py \\
python <lang>/tokenizer/train\_tokenizer.py --vocab-size 8000 \\
python scripts/validate\_phase1.py \\
python report/phase1/make\_plots.py
\end{quote}

Large artifacts (raw corpora, cleaned corpora, tokenizer files) are hosted
on OneDrive: \url{https://iiithydstudents-my.sharepoint.com/personal/swaraj_kumar_students_iiit_ac_in/_layouts/15/onedrive.aspx?id=%2Fpersonal%2Fswaraj%5Fkumar%5Fstudents%5Fiiit%5Fac%5Fin%2FDocuments%2FLMA%20DATA%20SET&ga=1}
(``LMA DATA SET'' shared OneDrive folder, requires an IIIT account with
access) -- see the repository root \texttt{README.md}'s artifact table for
current upload status.

% =====================================================================
\section{Limitations and What I Would Do Differently}
\label{sec:limitations}
% =====================================================================

\begin{itemize}
\item \textbf{The language-ID proxy is a real, acknowledged weak point.}
\texttt{langdetect} has no Bhojpuri label at all, so both directions of
this problem are handled by coarse Devanagari-script-ratio plus
hand-picked lexical markers, not a trained classifier. It is asymmetric
by design: Hindi's \texttt{likely\_bhojpuri} is an active reject filter,
Bhojpuri's \texttt{likely\_hindi\_contamination} is only a measurement.
\texttt{validate\_phase1.py}'s cross-language check caught 2 real misses
this heuristic let through -- there are almost certainly more that a
hash-based check cannot catch. Given more time, this would be the first
thing to replace with a real trained classifier.

\item \textbf{SentencePiece training samples the corpus}, most
aggressively for Hindi (22.0\% of train characters). This does not bias
the held-out evaluation numbers (measured separately, on unseen
validation text either way), but it does mean the trained model saw
under a quarter of Hindi's train split's characters.

\item \textbf{Near-duplicate detection beyond exact-hash matching was cut
for time.} Both languages only dedup on an exact normalized-text hash.
Given Bhojpuri's several overlapping \texttt{Satyam810/BhojpuriCorpus}
partitions and the BHLTR/jnu\_bhltr overlap, a MinHash/LSH pass would
likely find more duplication than the exact-hash pass did.

\item \textbf{The BPE-vs-Unigram argument (Section~\ref{sec:results}) is
structural, not empirically confirmed here.} No side-by-side BPE sweep
was run on this data to measure whether Unigram's advantage on
agglutinative suffixing and Devanagari conjuncts actually shows up as a
measurable difference for Hindi/Bhojpuri specifically.

\item \textbf{Two small raw-total discrepancies} (Hindi: 0.00003\%;
Bhojpuri: 0.024\%, both between the per-source breakdown sum and
\texttt{clean\_stats.json}'s own \texttt{chars\_in}) were not root-caused.
Both are too small to affect any pass/fail check.
\end{itemize}

% =====================================================================
\section{The Nepali Attempt}
\label{sec:nepali}
% =====================================================================

Alongside the graded Model L (Bhojpuri), a second, independent
low-resource candidate -- \textbf{Nepali}, also on the assignment's
allowed list -- was attempted using a more elaborate, config-driven
pipeline (\texttt{nepali/scripts/run\_pipeline.py}, stages S02--S14). This
was exploratory: Nepali has more public data available than Bhojpuri, so
it was worth checking whether it could clear the 20\% manual-share floor
with a stronger overall corpus than Bhojpuri's.

\textbf{What happened.} An early run of that pipeline measured only
\textbf{6.3\% manual share} -- failing the hard 20\% floor. This turned
out to be a \textbf{stale measurement}: most of the raw manual data (in
particular a 200{,}804-document news-scrape source, \texttt{onlinekhabar},
$\sim$72M words) had not finished collecting when that run happened. Once
$\sim$84.4M raw manual words were confirmed actually on disk (versus the
$\sim$15.4M that had fed the stale run), the pipeline was relaunched from
a clean state against the complete raw corpus.

\textbf{Progress on the rerun}, all confirmed from the pipeline's own
event log (\texttt{log.md}):

\begin{table}[htbp]
\centering
\begin{tabular}{lll}
\toprule
Stage & Result & Retention \\
\midrule
S02 Token calibration & Success & --- \\
S03 Normalize & Success & 1{,}658{,}065 / 1{,}658{,}065 docs, 0 dropped \\
S04 Language ID & Success & 99.77\% \\
S07 Holdout split & Success & train 1{,}406{,}033 / val 124{,}042 / test 124{,}221 \\
S05 Heuristic filters & Success & 99.72\% \\
S06 Perplexity filter & Success & 97.04\% \\
S08 Deduplication (MinHash/LSH) & \textbf{Stalled} & --- \\
\bottomrule
\end{tabular}
\caption{Nepali rerun progress, by pipeline stage.}
\end{table}

A real environment bug was hit and fixed along the way: \texttt{fasttext}
had no prebuilt wheel for this environment's Python 3.14, so the S04
language-ID stage alone was routed through a separate Python 3.11
install already on the machine -- the language-ID logic itself was not
bypassed, only its runtime. A second real bug was also fixed:
\texttt{nepali/configs/tokenizer.yaml} was configured to sweep a
16k--40k BPE vocabulary range, which violates the same course-TA ruling
covered in Section~\ref{sec:results}; it was corrected to
\texttt{[5000, 6000, 8000, 10000]} before the tokenizer stage could run
on it.

\textbf{Why it was not used for the final submission.} S08, the
MinHash/LSH deduplication pass over $\sim$1.36M documents, is
single-threaded for its bucket-insertion phase (needed to keep
manual-priority collision resolution deterministic) after a parallel
signature-computation phase. At its measured pace during the rerun, this
stage alone was on track to take on the order of 10+ hours to complete --
before stages S09 through S14 (decontamination, lane balancing, final
split, tokenizer training, tokenizer evaluation, corpus statistics) could
even begin. That was not a safe bet against the Phase 1 deadline.

\textbf{Bhojpuri is the final, graded Model L.} It was fully built,
cleaned, split, tokenized, and passing every
\texttt{scripts/validate\_phase1.py} check before the Nepali rerun was
even relaunched, and nothing in this report depends on Nepali finishing.
The Nepali pipeline's code (\texttt{nepali/scripts/}, \texttt{nepali/configs/})
is kept in the project repository for transparency into this exploration;
its data (raw and interim, $\sim$20GB) is not versioned, for the same
``code in git, corpus data elsewhere'' reason Hindi and Bhojpuri's own
raw/cleaned corpora are not versioned either.

% =====================================================================
\section*{Sources and References}
\addcontentsline{toc}{section}{Sources and References}
% =====================================================================

\begin{enumerate}
\item Hindi Wikipedia XML dump (self-extracted via WikiExtractor).
Source: Wikimedia dumps, \url{https://dumps.wikimedia.org}.
\item Bhojpuri Wikipedia XML dump (self-extracted via WikiExtractor).
Source: Wikimedia dumps, \url{https://dumps.wikimedia.org}.
\item Gadya Kosh (Hindi prose wiki). \url{https://gadyakosh.org}.
Author-copyright content, collected for coursework research use.
\item Kavita Kosh (Hindi/Urdu poetry wiki, also hosts Bhojpuri/Awadhi/
Maithili sections). \url{https://kavitakosh.org}. Author-copyright content,
collected for coursework research use.
\item Patrika (Hindi newspaper archive, Rajasthan Patrika group).
\url{https://www.patrika.com}. Collected for coursework research use.
\item Khabar Bhojpuri (Bhojpuri news site). \url{https://khabarbhojpuri.com}.
\item Anjoria (Bhojpuri content site). \url{https://anjoria.com}.
\item Jogira (Bhojpuri literature/culture site). \url{https://jogira.com}.
\item Sirijan (Bhojpuri quarterly e-magazine). \url{https://sirijan.com}.
\item Internet Archive (OCR/text extraction of books, journals, government
debate transcripts). \url{https://archive.org}.
\item HuggingFaceFW/fineweb-2 (Hindi and Bhojpuri partitions).
\url{https://huggingface.co/datasets/HuggingFaceFW/fineweb-2}.
\item AI4Bharat Sangraha (verified and unverified Hindi partitions).
\url{https://huggingface.co/datasets/ai4bharat/sangraha}.
\item Satyam810/BhojpuriCorpus (17 partitions: fish\_food, jnu\_bhltr,
hplt, wikipedia, eng\_bho, bhojpuri\_lr\_v3, kumar\_bhojpuri, iit\_bhu,
fleurs\_bhojpuri, google\_smol, glotcc, rural\_women, ud\_bhtb, and others).
\url{https://huggingface.co/datasets/Satyam810/BhojpuriCorpus}.
\item BHLTR: Bhojpuri-Hindi Language Technology Resources (JNU), monolingual
and parallel corpora. \url{https://github.com/shashwatup9k/BHLTR},
CC-BY-NC-SA-4.0.
\item \texttt{wikimedia/wikipedia:20231101.hi} -- evaluated, quarantined,
not used (redundant with the self-extracted Hindi Wikipedia dump above).
\url{https://huggingface.co/datasets/wikimedia/wikipedia}.
\item CulturaX, sil-ai/bloom-lm, AI4Bharat Sangraha (Bhojpuri rows),
Glot500, MADLAD-400, CC-100 -- evaluated for Bhojpuri and excluded; see
Section~\ref{sec:cleaning} for reasons per source.
\item SentencePiece (tokenizer library). Kudo, T. and Richardson, J.
(2018). \textit{SentencePiece: A simple and language independent subword
tokenizer and detokenizer for neural text processing.} EMNLP 2018 (System
Demonstrations). \url{https://github.com/google/sentencepiece}.
\item \texttt{langdetect} (Python language-identification library used as
one signal in the language filter, with known no-Bhojpuri-label
limitation documented throughout Section~\ref{sec:cleaning}).
\url{https://pypi.org/project/langdetect/}.
\end{enumerate}

\end{document}
